\subsection{Bonificación y salvaguardas}

Como se mostró en la Sección~4.5, una bonificación acotada puede restablecer el incentivo marginal una vez alcanzada la meta, sin modificar la construcción del ICG de resultados. Su objeto debe ser el resultado adicional efectivamente obtenido y no la ejecución del gasto, la adquisición de activos o el uso de recursos por sí mismos. Una inversión puede contribuir al sobrecumplimiento, pero el reconocimiento solo debería producirse cuando entre en operación y genere una mejora verificable.

El pago unitario puede tomar como referencia el costo regulatorio de la mejora programada, pero debería representar únicamente una fracción de dicho costo, de modo que parte del beneficio permanezca en favor de los usuarios. Asimismo, el premio debe incorporar un límite máximo que evite convertirlo en una fuente abierta de ingresos o inducir una concentración excesiva de recursos en los indicadores más fácilmente bonificables.

La selección de indicadores elegibles debería atender a su relevancia para los usuarios, controlabilidad y verificabilidad. No resulta conveniente premiar variables determinadas principalmente por factores externos ni resultados cuya medición no pueda sustentarse mediante información confiable y trazable. Cuando la evaluación dependa de equipos, muestras o registros operativos, las condiciones de recopilación, cobertura, frecuencia y conservación de la información deben definirse previamente.

El diseño también debe exigir que la mejora sea sostenible y que no se obtenga a costa del deterioro de otros indicadores. Para ello, el premio puede condicionarse al mantenimiento del sobrecumplimiento durante un periodo mínimo, al cumplimiento de estándares básicos en las demás dimensiones y a la ausencia de efectos adversos relevantes en otras localidades o componentes del servicio. La calibración del pago y la evaluación de las intervenciones también deben prevenir la sobreinversión y la utilización de activos innecesarios para alcanzar el resultado premiado.

Finalmente, la línea base, la meta y el tramo bonificable deberían determinarse mediante información histórica, comparaciones con empresas similares y criterios técnicos independientes del premio posterior, a fin de evitar una fijación estratégica de metas poco exigentes. Dadas estas exigencias, la implementación puede realizarse inicialmente mediante un programa piloto concentrado en pocos indicadores con alta relevancia, medición confiable y suficiente controlabilidad. La evaluación del piloto debería considerar el resultado adicional, su costo, sostenibilidad, distribución territorial y efectos sobre las demás dimensiones de la prestación.

En conjunto, la arquitectura recomendada corresponde al esquema \(B^{+}\): un ICG concentrado en resultados, un mecanismo separado y creíble de \textit{enforcement} para las inversiones, una tipificación congruente con la obligación incumplida y una bonificación acotada por mejoras adicionales verificables.
